\submittedby{Harisankar B}

We use the following characterisation of separations (Munkres, Lemma~23.1):
\emph{if $W$ is a subspace of $X$, then a separation of $W$ is a pair of disjoint non-empty
sets whose union is $W$, neither of which contains a limit point (in $X$) of the other.}
We also use that if $Y$ is connected and $W=C\cup D$ is a separation with $Y\subseteq W$, then
$Y$ lies entirely in $C$ or entirely in $D$.

\begin{proof}
We show that $Y\cup A$ is connected; the argument for $Y\cup B$ is identical with $A$ and $B$
interchanged.

Suppose, for contradiction, that $C,D$ form a separation of $Y\cup A$. Since $Y$ is connected,
$Y$ lies entirely in one of them; say $Y\subseteq C$. Then $D\subseteq(Y\cup A)\setminus Y
\subseteq A$. We claim that
\[ X=(C\cup B)\cup D \]
is a separation of $X$.

\emph{Union and disjointness.} $C\cup D=Y\cup A$ and $A\cup B=X\setminus Y$, so
$(C\cup B)\cup D=Y\cup A\cup B=X$. Moreover $D\cap C=\emptyset$, and
$D\cap B\subseteq A\cap B=\emptyset$.

\emph{Non-emptiness.} $D\neq\emptyset$ and $C\neq\emptyset$, as $C,D$ is a separation; hence
$C\cup B\neq\emptyset$.

\emph{No limit points across.} By Lemma~23.1 applied to the separations $C,D$ of $Y\cup A$ and
$A,B$ of $X\setminus Y$:
\begin{itemize}
  \item $D$ contains no limit point of $C$, and $C$ contains no limit point of $D$;
  \item $A$ contains no limit point of $B$, and $B$ contains no limit point of $A$.
\end{itemize}
Since $D\subseteq A$, the set $D$ contains no limit point of $B$; so $D$ contains no limit point
of $C\cup B$ (a limit point of a union is a limit point of one of the sets). Also every limit
point of $D$ is a limit point of $A$, so $B$ contains no limit point of $D$; together with the
first bullet, $C\cup B$ contains no limit point of $D$.

By Lemma~23.1, $C\cup B$ and $D$ form a separation of $X$, contradicting the connectedness of
$X$. Hence $Y\cup A$ is connected, and similarly $Y\cup B$ is connected.
\end{proof}
